\documentclass[11pt]{article}
\usepackage[margin=1in]{geometry}
\usepackage{url}
\usepackage{hyperref}

\title{A Minimal Execution Evidence Profile for Validator-Checkable AI Agent Operation Records in FAIR Digital Object-Inspired Data-Space Settings}
\author{Bin Zhang\\
Independent Researcher, China\\
Email: \href{mailto:joy7759@gmail.com}{joy7759@gmail.com}\\
ORCID: 0009-0002-8861-1481}
\date{}

\begin{document}
\maketitle

\begin{abstract}
The Execution Evidence and Operation Accountability Profile (EEOAP) is a
minimal evidence profile for validator-checkable records of individual
artificial intelligence (AI) agent operations in FAIR Digital Object-inspired
data-space settings. AI-agent outputs are often reviewed after the originating
runtime, logs, credentials, or prompts are unavailable. EEOAP packages actor,
action, subject, policy reference, output reference, provenance, evidence
references, integrity binding, and validation metadata into one offline review
object. The paper contributes the profile boundary, a local validator path,
and a repository-contained \texttt{paper\_case} artifact. The validator
enforces required structure, reference closure, and integrity-linked
consistency conditions over fields used by the review object. Running
\texttt{make paper-demo} demonstrates the intended boundary: the valid
evidence bundle reports \texttt{PASS valid evidence bundle}, while a
controlled tampered-output case reports \texttt{FAIL tampered output hash
mismatch} with \texttt{references\_digest\_mismatch} as the primary error
code. Targeted tests support the operation-accountability and paper-case
behavior, but exhaustive validator assurance, semantic correctness of AI
output, legal compliance, production readiness, public release, and archival
DOI claims are outside scope. The result is a standards- and
interfaces-oriented artifact for making selected post hoc operation evidence
claims inspectable and falsifiable offline.
\end{abstract}

\section{Introduction}

AI agent systems increasingly perform operations whose evidence matters after
the original interaction is over. An agent may read a data descriptor, call a
tool, apply a policy, produce a summary, export a review package, or hand a
result to another reviewer. A later reviewer may need to know which actor
performed the operation, what subject was acted on, which policy was
referenced, what output object was produced, which evidence references were
included, and whether the record still matches the referenced artifacts.

This paper introduces EEOAP (Execution Evidence and Operation Accountability
Profile), a minimal operation-level evidence profile and validator path for
independently reviewable AI agent operations in FAIR Digital Object-inspired
data-space settings. The profile records actor, action, subject, policy
reference, output, provenance, evidence references, artifacts, integrity
bindings, and validation metadata for one operation.

Runtime logs, traces, transcripts, provenance, schema validation,
attestation, and telemetry all provide useful background \cite{provdm,
provconstraints,jsonschema,slsa,slsaprov,intoto,otel,otelgenai}. The EEOAP
claim is narrower: a minimal evidence object can make selected
operation-level claims inspectable and falsifiable offline.

The current artifact fixes this claim to a reproducible paper case. Running
\texttt{make paper-demo} reports \texttt{PASS valid evidence bundle} for the
valid case and \texttt{FAIL tampered output hash mismatch} for the negative
case. The expected tampered failure code is
\texttt{references\_digest\_mismatch}. Targeted EEOAP tests previously passed
with 19 passed and 1 warning; full repository pytest success is not claimed.

\section{Problem and Scope}

The motivating problem is independent review of one AI agent operation after
execution. The reviewer may not have access to the original runtime
environment, live service credentials, private logs, raw prompts, or model
internals. A bounded record is needed to answer who or what acted, what action
was recorded, which subject was involved, which policy was referenced, what
output was produced, which evidence references belong to the record, and
whether the integrity binding still matches.

This paper treats that need as a profile and validation problem. A profile
narrows a broad space of execution data into a specified structure. A
validator checks whether a concrete object satisfies that structure, whether
references close, whether policy and evidence linkages are coherent, and
whether integrity values still match. The result is not a proof of truth,
legality, security, or semantic correctness.

EEOAP does not define organizational approval workflows, runtime enforcement,
human consent management, legal interpretation, agent reputation, policy
authoring, or production incident response. It does not claim conformance to,
certification by, endorsement from, or implementation of a specific FAIR
Digital Object Framework.

In this paper, ``validator-checkable'' refers only to offline checking of
structural completeness, reference closure, and selected linkage and integrity
constraints within the bounded EEOAP evidence object. It does not denote
formal verification of the underlying runtime, semantic correctness of the
generated output, legal compliance, or exhaustive assurance of the surrounding
system.

\section{EEOAP Profile Design}

EEOAP organizes execution evidence around one accountable operation. It is
designed to be small enough for review, strict enough for validation, and
explicit enough to separate claims from non-claims. The profile includes
actor, action, subject, policy reference, output reference, provenance,
evidence references, integrity binding, and validation metadata.

The term ``minimal'' is used in a pragmatic rather than formal proof-theoretic
sense. Each required element is included because it is needed either to
identify the acted-on object, reconstruct the claimed operation context,
connect the operation to a policy and evidence boundary, or recompute the
integrity boundary used in validation. In the paper case, removing actor,
subject, policy reference, output reference, evidence references, integrity
binding, or validation metadata would leave at least one review question
unanswered or would disable at least one validator check.

This profile is an interface boundary, not a complete representation of every
runtime event. An implementation may have richer logs, traces, prompts, tool
messages, approvals, and model telemetry. EEOAP selects a reviewable subset
and packages it as one evidence object.

The output reference is central to the tamper-detection example. The valid
evidence object includes an output digest that matches the operation output
content hash. The tampered evidence object changes the output digest while
leaving the integrity binding stale. The validator-backed path rejects the
object with \texttt{references\_digest\_mismatch}.

\section{Validator and Evidence Bundle}

The validator turns EEOAP from a descriptive document shape into a checkable
artifact. Within the paper-case boundary, the current validator path enforces
required structure, reference closure, and integrity-linked consistency
conditions over the fields used by the review object. The manuscript directly
demonstrates this boundary through a valid case and a controlled
integrity-failure case. Broader validator behavior is supported by targeted
tests but is not claimed exhaustively here.

The paper case is under \texttt{examples/paper\_case/}. It contains an
FAIR Digital Object-inspired subject descriptor, an aggregate-only policy
descriptor, an agent operation descriptor, a valid EEOAP evidence object, a
tampered object with a changed output digest and stale binding, and expected
validator pass/fail results.

The reproducible path is \texttt{make paper-demo}. The expected visible output
includes \texttt{PASS valid evidence bundle} and \texttt{FAIL tampered output
hash mismatch}. The expected primary error code is
\texttt{references\_digest\_mismatch}.

\section{FAIR Digital Object-inspired Data-space Mapping}

EEOAP is relevant to FAIR Digital Object-inspired data-space discussions
because those discussions often involve persistent object identity, typed
metadata, policy references, provenance, integrity, and exchange interfaces.
Digital object and object-interface work provides useful background
\cite{doip,kahnwilensky}.

In this manuscript, ``FAIR Digital Object-inspired'' denotes an object-centric
packaging pattern motivated by FAIR Digital Object discussions, including
persistent identity, machine-readable metadata, typed description, policy
references, and resolvable links. It does not claim conformance to a specific
FAIR Digital Object Framework implementation or to a DOIP (Digital Object
Interface Protocol, 数字对象接口协议) service.

The paper case uses \texttt{fdo-dataset.json} as the subject descriptor. It is
FAIR Digital Object-inspired rather than an implementation of a specific FAIR
Digital Object Framework. The descriptor gives the operation a subject
identity, type, owner, content hash, policy reference, and metadata. EEOAP
then records an agent operation against that subject and binds the result to
evidence references and integrity values.

This mapping does not claim FAIR Digital Object Framework conformance,
certification, endorsement, or deployment. It also does not implement DONA
DOIP.

\section{Reproducible Paper Case and Evaluation}

The evaluation is intentionally scoped to the current paper artifact. It asks
whether the repository contains a reproducible positive case and a
reproducible negative case for the core EEOAP claim. The evaluation is not a
model benchmark, user study, deployment study, or formal conformance test
against an external standard.

The verified evaluation facts are: \texttt{make paper-demo} passes; the valid
evidence bundle reports \texttt{PASS valid evidence bundle}; the tampered
output reports \texttt{FAIL tampered output hash mismatch}; the expected
tampered primary error code is \texttt{references\_digest\_mismatch}; and the
targeted EEOAP tests previously reported 19 passed and 1 warning. Full
repository pytest success is not claimed.

\section{Claim-to-Evidence Boundary}

EEOAP defines a minimal operation-level profile. The evidence is the
submission manuscript, the valid paper case evidence object, and associated
paper case files. The boundary is that EEOAP is not a full governance platform
or universal agent registry.

The validator enforces required structure, reference closure, and
integrity-linked consistency conditions for the paper case through the local
\texttt{make paper-demo} path. The boundary is that it does not prove legal
compliance, runtime security, semantic correctness, production readiness, or
exhaustive validator assurance.

The artifact has a local sealed anchor: tag \texttt{eeoap-v0.1-paper} at
commit \texttt{96f444b7ed39b39fe9f47e428af835952e843cb0}. The tag is not
claimed as publicly pushed, and no public GitHub Release or Zenodo DOI is
claimed.

\section{Related Work}

W3C PROV-DM and PROV-Constraints provide provenance and consistency-checking
background \cite{provdm,provconstraints}. JSON Schema Draft 2020-12 provides
JSON-style schema and validation background \cite{jsonschema}. ACM Artifact
Review and Badging provides artifact-reviewability and reproducibility
framing \cite{acmbadging}. DONA DOIP and Kahn and Wilensky provide digital
object and object-interface background \cite{doip,kahnwilensky}.

SLSA Version 1.2, SLSA Build Provenance, and in-toto Stable provide
supply-chain provenance and attestation background
\cite{slsa,slsaprov,intoto}. OpenTelemetry Semantic Conventions and the
OpenTelemetry generative AI conventions provide observability and telemetry
background \cite{otel,otelgenai}. EEOAP uses these works as related
background. It does not claim conformance to them.

EEOAP is adjacent to several existing approaches but narrower than all of
them. Unlike PROV, it does not attempt to represent arbitrary provenance
graphs or claim PROV validity; it packages one reviewable operation record.
Unlike JSON Schema, it does not rely on schema assertions alone, because
cross-reference and digest checks are part of the validator path. Unlike SLSA
and in-toto, it does not attest a software supply chain; it records post hoc
evidence for one AI-agent operation. Unlike OpenTelemetry, it is not a live
observability convention; it is an offline review object.

\section{Limitations and Threats to Validity}

The main limitation is scale. The current artifact contains one paper case,
not a large corpus of agent operations. It does not show generality across
agent frameworks, data domains, deployment models, or data-space
implementations.

The validator cannot prove that AI output is true, complete, unbiased,
policy-sufficient, or useful. It cannot prove that the runtime was secure,
that agent identity was uncompromised, that clocks were correct, that signing
keys were protected, or that evidence entered an external transparency
system.

ZKP is mentioned only as a non-claim and future direction. No ZKP
implementation is present. In broader deployments, operation evidence may
contain prompts, tool definitions, outputs, identifiers, or policy references
that carry sensitive or proprietary information. The present paper does not
solve selective disclosure, retention control, or privacy-preserving release
of such content and therefore treats privacy-preserving deployment as out of
scope. The FAIR Digital Object-inspired mapping is not FAIR Digital Object
Framework conformance, certification, endorsement, or deployment. Full
repository pytest success is not claimed.

\section{Standards and Interface Implications}

The standards-facing implication of EEOAP is that AI agent operations may need
a portable evidence object between raw runtime traces and broad governance
claims. For interface designers, the profile suggests that a completed
operation should expose at least actor, action, subject, policy reference,
output reference, provenance, evidence references, integrity binding, and
validation metadata.

For conformance and validation discussions, EEOAP suggests a staged approach:
define the smallest object that can be validated, provide a local checker that
accepts valid evidence and rejects a known tampered case, and add stronger
assurance layers only when the base object is stable.

\section{Conclusion}

EEOAP defines a minimal operation-level evidence profile for one AI agent
operation. It packages actor, action, subject, policy, output, provenance,
evidence references, integrity bindings, and validation metadata into an
object that can be reviewed outside the original runtime.

The reproducible paper case demonstrates the core boundary. Running
\texttt{make paper-demo} reports \texttt{PASS valid evidence bundle} for the
valid object and \texttt{FAIL tampered output hash mismatch} for the tampered
object. The expected primary error code is
\texttt{references\_digest\_mismatch}. Targeted EEOAP tests previously passed
with 19 passed and 1 warning; full repository pytest success is not claimed.

\section{Artifact Availability}

The artifact supporting the bounded claims of this paper consists of the
\texttt{paper\_case} files, the local validator path, and the
\texttt{make paper-demo} reproduction command. At initial submission, the
supporting software artifact is not deposited in a public repository because
the author is maintaining a sealed review-state package and does not yet
claim a public release or archival identifier. A private review package can be
supplied to editors and reviewers through the journal workflow on request.

This submission does not claim a public GitHub Release or a Zenodo DOI. The
claims supported by the artifact are limited to acceptance of the valid paper
case, rejection of a controlled tampered case with the expected error
boundary, and the bounded targeted-test evidence described in the manuscript.
If the artifact is later publicly released, this statement will be updated to
include the persistent access point.

\section*{Declaration of generative AI and AI-assisted technologies in the manuscript preparation process}

During the preparation of this work, the author used OpenAI ChatGPT and Codex
to support language editing, structural reorganization, command generation,
and drafting of submission-packaging text. After using these tools, the
author reviewed and edited the content as needed and takes full responsibility
for the content of the article.

\begin{thebibliography}{11}

\bibitem{provdm}
W3C (World Wide Web Consortium). PROV-DM: The PROV Data Model. W3C
Recommendation, 30 April 2013. Editors: Luc Moreau and Paolo Missier.
Available at: \url{https://www.w3.org/TR/prov-dm/}. Accessed 21 May 2026.

\bibitem{provconstraints}
W3C (World Wide Web Consortium). Constraints of the PROV Data Model. W3C
Recommendation, 30 April 2013. Editors: James Cheney, Paolo Missier, and Luc
Moreau; author: Tom De Nies. Available at:
\url{https://www.w3.org/TR/prov-constraints/}. Accessed 21 May 2026.

\bibitem{jsonschema}
Austin Wright, Henry Andrews, Ben Hutton, and Greg Dennis. JSON Schema Draft
2020-12. JSON Schema project, published 16 June 2022. Available at:
\url{https://json-schema.org/draft/2020-12}. Accessed 21 May 2026.

\bibitem{acmbadging}
ACM (Association for Computing Machinery). Artifact Review and Badging -
Current. Artifact Review and Badging Version 1.1, 24 August 2020. Available
at: \url{https://www.acm.org/publications/policies/artifact-review-and-badging-current}.
Accessed 21 May 2026.

\bibitem{doip}
DONA Foundation. Digital Object Interface Protocol Specification. Version 2.0,
12 November 2018. Available at:
\url{https://www.dona.net/sites/default/files/2018-11/DOIPv2Spec_1.pdf}.
Accessed 21 May 2026.

\bibitem{kahnwilensky}
Robert Kahn and Robert Wilensky. A framework for distributed digital object
services. International Journal on Digital Libraries, 6, 115--123, 2006.
\url{https://doi.org/10.1007/s00799-005-0128-x}.

\bibitem{slsa}
SLSA (Supply-chain Levels for Software Artifacts). SLSA specification. Version
1.2, status: Approved. Available at: \url{https://slsa.dev/spec/v1.2/}.
Accessed 21 May 2026.

\bibitem{slsaprov}
SLSA (Supply-chain Levels for Software Artifacts). Build: Provenance. Status:
Approved. Predicate type: \url{https://slsa.dev/provenance/v1}. Available at:
\url{https://slsa.dev/spec/v1.2/build-provenance}. Accessed 21 May 2026.

\bibitem{intoto}
in-toto. Specifications: in-toto Stable (v1.0). Status from source:
thoroughly reviewed specification. Available at:
\url{https://in-toto.io/docs/specs/}. Accessed 21 May 2026.

\bibitem{otel}
OpenTelemetry. OpenTelemetry semantic conventions 1.41.0. Available at:
\url{https://opentelemetry.io/docs/specs/semconv/}. Accessed 21 May 2026.

\bibitem{otelgenai}
OpenTelemetry. Semantic conventions for generative AI systems. OpenTelemetry
semantic conventions 1.41.0, status: Development. Available at:
\url{https://opentelemetry.io/docs/specs/semconv/gen-ai/}. Accessed 21 May
2026.

\end{thebibliography}

\end{document}
